Fix \(P\in\mathcal P_L\), and abbreviate \(H=H_{n,\Delta}\), \(p=p_{n,\Delta}=H^2/2\), and \(B=2LH/3\).  Since \(H\) is a quantizer boundary, \(\{Q_\Delta(R)<H\}=\{R<H\}\).  Thus \(\mathrm{lem:radial\text{-}grouped\text{-}law}\) and the iid and Lipschitz member properties in \(\mathrm{def:half\text{-}disc\text{-}class}\) give, for each \(t\),
\[
 S_{t,n,\Delta}\sim\operatorname{Binomial}(n,q_t),
 \qquad
 q_t=p\mu_t,
 \qquad
 \mu_t=\frac1p\int_0^H r m_t(r)\,dr,
 \qquad
 |\mu_t-m_t(0)|\le B.
\]
Since \(p>0\), all divisions in \(\mathrm{def:total\text{-}honest\text{-}interval}\) are valid.  The Clopper--Pearson interval is defined at every count, including zero and \(n\).  Its defining exact-binomial property gives
\[
 \mathbb P_P\{q_t\in[\ell_t,u_t]\}\ge1-\frac\alpha2.
\]
No independence between the two side counts is required: the union bound gives simultaneous inclusion of \(q_0,q_1\) with probability at least \(1-\alpha\).  On that event,
\[
 \mu_1-\mu_0\in
 [\ell_1/p-u_0/p,\ u_1/p-\ell_0/p]
\]
and
\[
 |\tau_P(b)-(\mu_1-\mu_0)|
 \le |m_1(0)-\mu_1|+|m_0(0)-\mu_0|
 \le2B.
\]
Thus \(\tau_P(b)\in C_{n,\Delta,\alpha}\) with probability at least \(1-\alpha\).  The endpoints are measurable functions of the released counts and fixed constants, so the interval is a closed bounded, \(\sigma(Z_\Delta^n)\)-measurable interval for every sample.  This proves honesty uniformly over \(\mathcal P_L\).

For completeness, we derive the needed sparse-bin length bound rather than invoking an asymptotic approximation.  Let \(S\sim\operatorname{Binomial}(n,q)\), write \(\widehat q=S/n\), let \([\ell,u]=I^{\mathrm{CP}}_{n,\gamma}(S)\), and put \(\beta=\gamma/2\).  For \(S<n\), the upper endpoint satisfies
\[
 \mathbb P_u\{W\le S\}=\beta,
 \qquad W\sim\operatorname{Binomial}(n,u),
\]
and \(u>\widehat q\).  On \(W\le S\), \(|W-nu|\ge n(u-\widehat q)\); hence the pointwise bound for the indicator of this event and \(\mathbb E_u(W-nu)^2=nu(1-u)\) imply
\[
 \beta\le\frac{nu(1-u)}{n^2(u-\widehat q)^2}.
\]
Consequently, with \(d=u-\widehat q\),
\[
 d^2\le\frac{\widehat q+d}{n\beta},
 \qquad
 d\le\sqrt{\frac{\widehat q}{n\beta}}+\frac1{n\beta}.
\]
The same argument at the lower endpoint, for \(S>0\), uses \(\mathbb P_\ell\{W\ge S\}=\beta\) and gives
\[
 (\widehat q-\ell)^2\le\frac{\ell(1-\ell)}{n\beta}
 \le\frac{\widehat q}{n\beta}.
\]
The boundary conventions \(\ell=0\) at \(S=0\) and \(u=1\) at \(S=n\) obey the same resulting inequality.  Therefore, for every count,
\[
 u-\ell\le2\sqrt{\frac{\widehat q}{n\beta}}+\frac1{n\beta}.
\]
Taking expectations and using concavity of the square root gives
\[
 \mathbb E_q(u-\ell)
 \le2\sqrt{\frac{q}{n\beta}}+\frac1{n\beta}.
\]
Apply this with \(\gamma=\alpha/2\), so \(\beta=\alpha/4\), and with \(q=q_t\le p\).  The interval length is exactly
\[
 \operatorname{length}(C_{n,\Delta,\alpha})
 =\frac{(u_1-\ell_1)+(u_0-\ell_0)}p+4B.
\]
It follows that
\[
 \mathbb E_P\operatorname{length}(C_{n,\Delta,\alpha})
 \le\frac4{\sqrt{n\beta p}}+\frac2{n\beta p}+4B.
\]
Finally, \(h_n\le H\le h_n+\Delta\le2\rho_{n,\Delta}\), \(p=H^2/2\), and \(h_n=n^{-1/4}\) imply
\[
 \frac1{\sqrt n H}\le h_n,
 \qquad
 \frac1{nH^2}\le h_n^2\le\rho_{n,\Delta}.
\]
Thus
\[
 \sup_{P\in\mathcal P_L}\mathbb E_P\operatorname{length}(C_{n,\Delta,\alpha})
 \le\left(4\sqrt{\frac2\beta}+\frac4\beta+\frac{16L}{3}\right)\rho_{n,\Delta},
 \qquad \beta=\frac\alpha4.
\]
This is finite in every \(\Delta\)-regime and proves the claimed expected-length bound.